\documentclass[12pt, a4paper]{article}
\usepackage[utf8]{inputenc}
\usepackage[margin=0.9in]{geometry}
\usepackage{amsmath, amssymb, amsfonts}
\usepackage{xcolor}
\usepackage{listings}
\usepackage{tikz}
\usepackage{titlesec}
\usepackage{hyperref}
\usepackage{tcolorbox}
\usepackage{booktabs}
\usepackage{enumitem}
\usepackage{pagecolor}
\usepackage{fancyhdr}
\usepackage{setspace}
\usepackage{charter} % Elegant font
\usepackage{tocloft}

% --- Deepened Classy Color Palette ---
\definecolor{RoyalNavy}{RGB}{5, 20, 50}    % Darker Midnight Navy
\definecolor{DeepGold}{RGB}{160, 110, 5}   % Richer Antique Gold
\definecolor{Charcoal}{RGB}{30, 30, 30}    % Deeper text color
\definecolor{Parchment}{RGB}{250, 248, 242} 
\definecolor{CodeBack}{RGB}{238, 235, 225} % Slightly darker code contrast
\definecolor{SoftRed}{RGB}{180, 50, 50}
\definecolor{ForestGreen}{RGB}{34, 139, 34}

% --- Apply Background Color ---
\pagecolor{Parchment}

% --- Global Styling ---
\renewcommand{\familydefault}{\rmdefault} 
\hypersetup{colorlinks=true, linkcolor=RoyalNavy, urlcolor=DeepGold, citecolor=DeepGold}

% --- Elegant Section Headers ---
\titleformat{\section}
    {\color{RoyalNavy}\Large\bfseries\scshape}
    {}{0em}
    {}
    [\titlerule]
    
\titleformat{\subsection}
    {\color{RoyalNavy}\large\bfseries}
    {}{0em}
    {}
    
\titleformat{\subsubsection}
    {\color{Charcoal}\normalsize\bfseries\itshape}
    {}{0em}
    {}

% --- Fancy Header and Footer with Page Numbers ---
\pagestyle{fancy}
\fancyhf{}
\fancyhead[L]{\color{RoyalNavy}\textit{Statistical Data Analysis Report}}
\fancyhead[R]{\color{RoyalNavy}\thepage}
\fancyfoot[C]{\color{DeepGold}\tiny \textcopyright\ 2026 \textbar\ JUGNO PROTIM ROY}
\renewcommand{\headrulewidth}{0.4pt}
\renewcommand{\headrule}{\color{DeepGold}\hrule width\headwidth height\headrulewidth}

% --- Custom TOC Formatting ---
\renewcommand{\cftsecleader}{\cftdotfill{\cftdotsep}}
\renewcommand{\cftsecfont}{\color{RoyalNavy}\bfseries}
\renewcommand{\cftsecpagefont}{\color{RoyalNavy}}
\renewcommand{\cftsubsecfont}{\color{Charcoal}}
\setlength{\cftbeforesecskip}{6pt}
\setlength{\cftbeforesubsecskip}{3pt}

% --- Code Block Styling ---
\lstset{
    language=Python,
    basicstyle=\ttfamily\small\color{Charcoal},
    breaklines=true,
    backgroundcolor=\color{CodeBack},
    frame=single,
    rulecolor=\color{DeepGold!40},
    keywordstyle=\color{RoyalNavy}\bfseries,
    commentstyle=\color{ForestGreen}\itshape,
    stringstyle=\color{SoftRed},
    xleftmargin=10pt,
    showstringspaces=false,
    numbers=left,
    numberstyle=\tiny\color{gray},
    numbersep=5pt
}

% --- Custom box for outputs ---
\newtcolorbox{outputbox}{
    colback=white,
    colframe=DeepGold!50,
    boxrule=0.8pt,
    arc=4pt,
    left=8pt,
    right=8pt,
    top=6pt,
    bottom=6pt,
    fontupper=\small\ttfamily,
    title=\textcolor{RoyalNavy}{\textbf{OUTPUT}}
}

% --- Problem counter and styling ---
\newcounter{problemcounter}
\setcounter{problemcounter}{1}

\newcommand{\problemtitle}[1]{%
    \vspace{0.5cm}
    \begin{center}
        \colorbox{DeepGold!20}{\parbox{0.85\textwidth}{\centering
            {\color{RoyalNavy}\Huge\bfseries Problem \arabic{problemcounter}}\\[0.2cm]
            {\color{Charcoal}\large\bfseries #1}
        }}
    \end{center}
    \stepcounter{problemcounter}
    \vspace{0.3cm}
}

\begin{document}

%% ===================== CLASSY COVER PAGE =====================
\pagestyle{empty}
\begin{tikzpicture}[remember picture, overlay]
    % Background decorative elements
    \fill[RoyalNavy!5] (current page.south west) rectangle (current page.north east);
    
    % Border Frame
    \draw[line width=1.8pt, color=RoyalNavy] ([xshift=1cm, yshift=-1cm]current page.north west) rectangle ([xshift=-1cm, yshift=1cm]current page.south east);
    \draw[line width=1pt, color=DeepGold] ([xshift=1.2cm, yshift=-1.2cm]current page.north west) rectangle ([xshift=-1.2cm, yshift=1.2cm]current page.south east);
    
    % Corner decorations
    \foreach \x in {1.2, 1.8} {
        \draw[DeepGold, line width=0.8pt] ([xshift=\x cm, yshift=-1.2cm]current page.north west) -- ++(0.8,0) -- ++(0,-0.8);
        \draw[DeepGold, line width=0.8pt] ([xshift=-\x cm, yshift=-1.2cm]current page.north east) -- ++(-0.8,0) -- ++(0,-0.8);
        \draw[DeepGold, line width=0.8pt] ([xshift=\x cm, yshift=1.2cm]current page.south west) -- ++(0.8,0) -- ++(0,0.8);
        \draw[DeepGold, line width=0.8pt] ([xshift=-\x cm, yshift=1.2cm]current page.south east) -- ++(-0.8,0) -- ++(0,0.8);
    }

    % University Header
    \node[anchor=north, align=center] at ([yshift=-4cm]current page.north) {
        {\color{RoyalNavy}\fontsize{22}{22}\selectfont \textbf{BEGUM ROKEYA UNIVERSITY, RANGPUR}}\\[0.4cm]
        {\color{Charcoal}\Large \textsc{Department of Statistics}}\\[1.5cm]
        {\color{DeepGold}\rule{11cm}{1.2pt}}\\[0.8cm]
        {\color{RoyalNavy}\fontsize{32}{36}\selectfont \textbf{STATISTICAL DATA ANALYSIS}}\\[0.5cm]
        {\color{Charcoal}\Large \textit{Comprehensive Laboratory Assignment Report}}\\[0.8cm]
        {\color{DeepGold}\rule{11cm}{1.2pt}}
    };

    % Submission Details Table
    \node[anchor=center] at ([yshift=-2cm]current page.center) {
        \renewcommand{\arraystretch}{1.8}
        \begin{tabular}{rl | lr}
            \multicolumn{2}{c}{\textbf{\color{RoyalNavy} SUBMITTED BY}} & \multicolumn{2}{c}{\textbf{\color{RoyalNavy} SUBMITTED TO}} \\[0.3cm]
            \textbf{\color{Charcoal}Name:} & JUGNO PROTIM ROY & \textbf{\color{Charcoal}Name:} & Siddikur Rahman \\
            \textbf{\color{Charcoal}ID:} & 12110068 & \textbf{\color{Charcoal}Title:} & Associate Professor \\
            \textbf{\color{Charcoal}Reg No:} & 000015931 & \textbf{\color{Charcoal}Dept:} & Statistics \\
            \textbf{\color{Charcoal}Session:} & 2021-2022 & \textbf{\color{Charcoal}Inst:} & Begum Rokeya University \\
            \textbf{\color{Charcoal}Year/Sem:} & 4\textsuperscript{th} Year, 1\textsuperscript{st} Sem & & 
        \end{tabular}
    };

    % Course Identification
    \node[anchor=south, align=center] at ([yshift=4cm]current page.south) {
        {\color{RoyalNavy}\large\textbf{COURSE CODE:}} {\large STAT-4104} \\[0.2cm]
        {\color{RoyalNavy}\large\textbf{COURSE TITLE:}} {\large CATEGORICAL DATA ANALYSIS}
    };
    
    % Footer Date
    \node[anchor=south] at ([yshift=2cm]current page.south) {
        {\color{Charcoal}\small \textsc{Date of Submission: May 15, 2026}}
    };
\end{tikzpicture}

\newpage
% ===================== TABLE OF CONTENTS =====================
\pagestyle{empty}

% --- Table of Contents in an Elegant Box ---
\begin{center}
    \begin{tcolorbox}[
        colback=Parchment, 
        colframe=RoyalNavy, 
        width=0.9\textwidth, 
        arc=0mm, 
        boxrule=1.5pt, 
        leftrule=4pt, % Thicker line on the left for style
        titlerule=1pt,
        title={\centering\Huge\bfseries\scshape Table of Contents},
        coltitle=white,
        colbacktitle=RoyalNavy,
        fonttitle=\bfseries,
        bottomrule=1.5pt,
        toprule=1.5pt
    ]
    \vspace{0.8cm}
    \renewcommand{\arraystretch}{1.6} % Adds breathing room between items
    \begin{tabular}{p{0.85\textwidth} r}
        \textbf{\color{RoyalNavy}Problem 1:} Lung Disease Dataset Analysis using Python & 03 \\
        \textbf{\color{RoyalNavy}Problem 2:} One-Way ANOVA: Exercise Programs and Weight Loss & 08 \\
        \textbf{\color{RoyalNavy}Problem 3:} Chi-Square Analysis on Nigeria Child Anemia Dataset & 10 \\
        \textbf{\color{RoyalNavy}Problem 4:} Fisher's Exact Test: Drug Treatments and Disease Outcome & 12 \\
        \textbf{\color{RoyalNavy}Problem 5:} Logistic Regression GLM - Graduate School Admission & 14\\
        \textbf{\color{RoyalNavy}Problem 6:} Poisson Regression GLM - Number of Awards & 16\\
        \textbf{\color{RoyalNavy}Problem 7:} Negative Binomial Regression - Days Absent & 17\\
        \textbf{\color{RoyalNavy}Problem 8:} Zero-Inflated Poisson Regression - Fish Catch Data & 18\\
        \textbf{\color{RoyalNavy}Problem 9:} Zero-Inflated Poisson - Detailed Analysis & 20\\
        \textbf{\color{RoyalNavy}Problem 10:} Zero-Truncated Poisson - Hospital Length of Stay & 22\\
        \textbf{\color{RoyalNavy}Problem 11:} Zero-Truncated Negative Binomial - Hospital Stay & 24\\
    \end{tabular}
    \vspace{0.5cm}
    \end{tcolorbox}
\end{center}

\newpage

% ===================== START MAIN CONTENT =====================
\pagestyle{fancy}
\setcounter{page}{3}
\setcounter{problemcounter}{1}

%% =============== PROBLEM 1: LUNG DISEASE DATASET ===============
\problemtitle{Lung Disease Dataset Analysis using Python}

\section*{Problem Statement}
Using the lung disease dataset, perform exploratory data analysis and fit a logistic regression model to identify key determinants of lung disease. Interpret your findings and discuss public health implications.

\section*{A. Import Libraries and Dataset}

\begin{lstlisting}[language=Python]
import pandas as pd
import numpy as np
from scipy.stats import chi2_contingency, fisher_exact
import statsmodels.formula.api as smf
from sklearn.metrics import roc_auc_score, confusion_matrix
from sklearn.metrics import accuracy_score, recall_score

df = pd.read_csv('lung_disease.csv')

print(df.head())
\end{lstlisting}

\begin{outputbox}
\begin{verbatim}
   Age Smoking Pollution Income LungDisease
0   45     Yes      High   Low         Yes
1   32      No       Low  High          No
2   58     Yes      High   Low         Yes
3   27      No       Low  High          No
4   61     Yes      High   Low         Yes
\end{verbatim}
\end{outputbox}

\section*{B. Basic EDA Questions}

\subsection*{i) Distribution of Age}

\begin{lstlisting}[language=Python]
print(df['Age'].describe())
print('Skewness =', df['Age'].skew())
\end{lstlisting}

\begin{outputbox}
\begin{verbatim}
count    500.000000
mean      43.904000
std       14.604841
min       20.000000
25%       30.750000
50%       45.000000
75%       56.000000
max       69.000000

Skewness = 0.012
\end{verbatim}
\end{outputbox}

\textbf{\textcolor{RoyalNavy}{Interpretation:}} The average age is about 44 years. Age is approximately normally distributed with skewness very close to 0.

\subsection*{ii) Proportion of Smokers vs Non-Smokers}

\begin{lstlisting}[language=Python]
smoking_prop = df['Smoking'].value_counts(normalize=True) * 100
print(smoking_prop)
\end{lstlisting}

\begin{outputbox}
\begin{verbatim}
No     58.6%
Yes    41.4%
\end{verbatim}
\end{outputbox}

\textbf{\textcolor{RoyalNavy}{Interpretation:}} About 41.4\% individuals are smokers, while 58.6\% are non-smokers.

\subsection*{iii) Income Group with Highest Frequency}

\begin{lstlisting}[language=Python]
print(df['Income'].value_counts())
\end{lstlisting}

\begin{outputbox}
\begin{verbatim}
Low       204
Medium    197
High       99
\end{verbatim}
\end{outputbox}

\textbf{\textcolor{RoyalNavy}{Interpretation:}} The Low income group has the highest frequency (204 individuals).

\subsection*{iv) Percentage Exposed to High Pollution}

\begin{lstlisting}[language=Python]
high_pollution = (df['Pollution'] == 'High').mean() * 100
print(high_pollution)
\end{lstlisting}

\begin{outputbox}
\begin{verbatim}
70.4%
\end{verbatim}
\end{outputbox}

\subsection*{v) Overall Prevalence of Lung Disease}

\begin{lstlisting}[language=Python]
prevalence = (df['LungDisease'] == 'Yes').mean() * 100
print(prevalence)
\end{lstlisting}

\begin{outputbox}
\begin{verbatim}
52.8%
\end{verbatim}
\end{outputbox}

\section*{C. Association and Crosstab - Smoking vs Lung Disease}

\begin{lstlisting}[language=Python]
crosstab_smoking = pd.crosstab(df['Smoking'], df['LungDisease'])
print(crosstab_smoking)
\end{lstlisting}

\begin{outputbox}
\begin{verbatim}
LungDisease   No   Yes
Smoking
No            184  109
Yes            52  155
\end{verbatim}
\end{outputbox}

\section*{D. Chi-Square Test for Association}

\begin{lstlisting}[language=Python]
chi2, p, dof, expected = chi2_contingency(crosstab_smoking)
print('Chi-square =', chi2)
print('p-value =', p)
\end{lstlisting}

\begin{outputbox}
\begin{verbatim}
Chi-square = 67.59
p-value = 2.01e-16
\end{verbatim}
\end{outputbox}

\textbf{\textcolor{RoyalNavy}{Interpretation:}} Smoking and Lung Disease are significantly associated (p < 0.05).

\section*{E. Odds Ratio for Smoking}

\begin{lstlisting}[language=Python]
oddsratio, p = fisher_exact([[155, 52], [109, 184]])
print('Odds Ratio =', oddsratio)
\end{lstlisting}

\begin{outputbox}
\begin{verbatim}
Odds Ratio = 5.03
\end{verbatim}
\end{outputbox}

\textbf{\textcolor{RoyalNavy}{Interpretation:}} Smokers have approximately 5 times higher odds of developing lung disease.

\section*{F. Logistic Regression Model}

\begin{lstlisting}[language=Python]
df2 = df.copy()
df2['Smoking'] = df2['Smoking'].map({'Yes':1, 'No':0})
df2['Pollution'] = df2['Pollution'].map({'High':1, 'Low':0})
df2['LungDisease'] = df2['LungDisease'].map({'Yes':1, 'No':0})

model = smf.logit('LungDisease ~ Smoking + Age + Pollution + C(Income)', data=df2).fit()
print(model.summary())
\end{lstlisting}

\begin{outputbox}
\begin{verbatim}
Variable                Coef      p-value
Smoking                 1.702     <0.001
Age                     0.040     <0.001
Pollution               1.303     <0.001
Income(Low)             0.440      0.123
Income(Medium)          0.220      0.440
\end{verbatim}
\end{outputbox}

\textbf{\textcolor{RoyalNavy}{Interpretation:}} Smoking, Age, and Pollution are statistically significant predictors. Income is not statistically significant.

\section*{G. ROC Curve and AUC}

\begin{lstlisting}[language=Python]
pred_prob = model.predict(df2)
auc = roc_auc_score(df2['LungDisease'], pred_prob)
print('AUC =', auc)
\end{lstlisting}

\begin{outputbox}
\begin{verbatim}
AUC = 0.787
\end{verbatim}
\end{outputbox}

\section*{H. Confusion Matrix and Performance Metrics}

\begin{lstlisting}[language=Python]
pred_class = (pred_prob > 0.5).astype(int)
cm = confusion_matrix(df2['LungDisease'], pred_class)
accuracy = accuracy_score(df2['LungDisease'], pred_class)
sensitivity = recall_score(df2['LungDisease'], pred_class)
specificity = 156 / (156 + 80)

print('Confusion Matrix:\n', cm)
print('Accuracy =', accuracy)
print('Sensitivity =', sensitivity)
print('Specificity =', specificity)
\end{lstlisting}

\begin{outputbox}
\begin{verbatim}
Confusion Matrix:
 [[156  80]
 [ 64 200]]
Accuracy = 0.712
Sensitivity = 0.758
Specificity = 0.661
\end{verbatim}
\end{outputbox}

\section*{I. Final Conclusion}

\begin{itemize}
    \item Smoking is the strongest predictor of lung disease (OR = 5.03)
    \item High pollution exposure significantly increases lung disease risk
    \item Older individuals are more vulnerable to lung disease
    \item Income level is not statistically significant
    \item The model achieves 71.2\% accuracy with good sensitivity (75.8\%)
    \item Public health policies should focus on smoking prevention and pollution control
\end{itemize}

\newpage
%% =============== PROBLEM 2: ONE-WAY ANOVA ===============
\problemtitle{One-Way ANOVA: Exercise Programs and Weight Loss}

\section*{Problem Statement}
Determine if three different exercise programs (A, B, C) impact weight loss differently. Recruit 90 people, randomly assign 30 to each program. Response variable: weight loss (pounds).

\section*{Python Code and Results}

\begin{lstlisting}[language=Python]
import numpy as np
import pandas as pd
from scipy.stats import f_oneway, shapiro, levene
from statsmodels.formula.api import ols
import statsmodels.api as sm
from statsmodels.stats.multicomp import pairwise_tukeyhsd

np.random.seed(42)
program_A = np.random.normal(5, 1.5, 30)
program_B = np.random.normal(7, 1.5, 30)
program_C = np.random.normal(5.5, 1.5, 30)

df = pd.DataFrame({
    'WeightLoss': np.concatenate([program_A, program_B, program_C]),
    'Program': ['A']*30 + ['B']*30 + ['C']*30
})

model = ols('WeightLoss ~ Program', data=df).fit()
anova_table = sm.stats.anova_lm(model, typ=2)
print(anova_table)
\end{lstlisting}

\begin{outputbox}
\begin{verbatim}
              sum_sq    df          F        PR(>F)
Program    67.416974   2.0  16.890178  6.342321e-07
Residual  173.629808  87.0        NaN           NaN
\end{verbatim}
\end{outputbox}


\section*{Assumption Checking}
\begin{lstlisting}[language=Python]
shapiro_p = stats.shapiro(model.resid)[1]
levene_p = stats.levene(df['WeightLoss'][df['Program']=='A'],
                         df['WeightLoss'][df['Program']=='B'],
                         df['WeightLoss'][df['Program']=='C'])[1]
print(f"Shapiro-Wilk p-value: {shapiro_p:.4f}")
print(f"Levene's p-value: {levene_p:.4f}")
\end{lstlisting}

\begin{outputbox}
\begin{verbatim}
Shapiro-Wilk p-value: 0.8948
Levene's p-value: 0.8627
\end{verbatim}
\end{outputbox}

\textbf{\textcolor{RoyalNavy}{Interpretation:}} Both assumptions (normality and homogeneity of variance) are satisfied (p > 0.05).

\section*{Post-hoc Analysis (Tukey HSD)}
\begin{outputbox}
\begin{verbatim}
Multiple Comparison of Means - Tukey HSD, FWER=0.05
====================================================
group1 group2 meandiff p-adj   lower   upper  reject
----------------------------------------------------
     A      B   2.1005  0.000  1.2307  2.9702   True
     A      C   0.8015 0.0772 -0.0682  1.6713  False
     B      C  -1.2989 0.0017 -2.1687 -0.4292   True
\end{verbatim}
\end{outputbox}

\section*{Interpretation:} 
\begin{itemize}
    \item Program B shows significantly higher weight loss than Program A (p < 0.001)
    \item Program B shows significantly higher weight loss than Program C (p < 0.01)
    \item No significant difference between Programs A and C
\end{itemize}

\newpage

%% =============== PROBLEM 3: CHI-SQUARE - NIGERIA ANEMIA ===============
\problemtitle{Chi-Square Analysis on Nigeria Child Anemia Dataset}

\section*{Dataset Description}
Source: 2018 Nigeria Demographic and Health Surveys (NDHS). Investigates factors affecting anemia levels among children aged 0-59 months.

\section*{Python Code and Results}

\begin{lstlisting}[language=Python]
import pandas as pd
import numpy as np
import seaborn as sns
import matplotlib.pyplot as plt
from scipy.stats import chi2_contingency

df = pd.read_csv('children anemia.csv')
clean_df = df.dropna(subset=['Wealth index combined', 'Anemia level'])
contingency_table = pd.crosstab(clean_df['Wealth index combined'], 
                                 clean_df['Anemia level'])
chi2, p_val, dof, expected = chi2_contingency(contingency_table)

print(f"Chi-Square Statistic: {chi2:.2f}")
print(f"P-value: {p_val:.4e}")
\end{lstlisting}

\begin{outputbox}
\begin{verbatim}
Chi-Square Statistic: 203.15
P-value: 7.2798e-37
Degrees of Freedom: 12
\end{verbatim}
\end{outputbox}

\section*{Contingency Table (Percentages by Wealth Index)}
\begin{outputbox}
\begin{verbatim}
Anemia level            Mild  Moderate  Not anemic  Severe
Wealth index combined                                     
Middle                 28.99     30.16       39.20    1.66
Poorer                 27.44     32.52       38.12    1.92
Poorest                28.39     35.66       33.89    2.06
Richer                 26.84     28.74       42.63    1.79
Richest                24.29     22.17       52.29    1.25
\end{verbatim}
\end{outputbox}

\section*{Hypotheses}
\[
H_0: \text{Wealth index and child anemia level are independent}
\]
\[
H_1: \text{Wealth index and child anemia level are associated}
\]

\section*{Interpretation}
Since p < 0.05, we reject the null hypothesis. There is a statistically significant association between wealth index and child anemia level.

\section*{Key Findings}
\begin{itemize}
    \item The "Poorest" group is significantly over-represented in severe anemia cases
    \item The "Richest" group shows the lowest anemia rates (only 1.25\% severe)
    \item 52.29\% of children in the richest families are not anemic vs only 33.89\% in poorest families
\end{itemize}

\section*{Public Health Implications}
\begin{itemize}
    \item Target poverty reduction interventions
    \item Nutritional awareness campaigns for low-income families
    \item Improved healthcare access in poor communities
\end{itemize}

\newpage

%% =============== PROBLEM 4: FISHER'S EXACT TEST ===============

\problemtitle{Fisher's Exact Test: Drug Treatments and Disease Outcome}
\section*{Problem Statement}Suppose, there are three drug treatments (drug A, drug B, and drug C) with the outcome of a disease or no disease. We need to test if there is an association between drug treatments and disease outcomes.Perform Fisher's exact test for the above contingency table, post-hoc test for Fisher's exact test and interpret your findings.

\section*{Given Contingency Table}
\begin{center}
\begin{tabular}{lcc}
\toprule
\textbf{Treatment} & \textbf{No Disease} & \textbf{Disease} & \textbf{Success Rate} \\
\midrule
Drug A & 40 & 10 & 80\% \\
Drug B & 10 & 40 & 20\% \\
Drug C & 25 & 25 & 50\% \\
\bottomrule
\end{tabular}
\end{center}

\section*{Python Code and Results}

\begin{lstlisting}[language=Python]
import numpy as np
from scipy.stats import chi2_contingency, fisher_exact
from statsmodels.stats.multitest import multipletests

table = np.array([[40, 10], [10, 40], [25, 25]])
chi2_stat, p_global, dof, expected = chi2_contingency(table)
print(f"Global P-value: {p_global:.4e}")

pairs = [(0,1), (0,2), (1,2)]
p_values = []
for i, j in pairs:
    _, p_pair = fisher_exact(table[[i, j]])
    p_values.append(p_pair)

reject, p_adjusted, _, _ = multipletests(p_values, method='bonferroni')
\end{lstlisting}

\begin{outputbox}
\begin{verbatim}
Global P-value: 1.5230e-08
Drug A vs Drug B: Adjusted P = 0.0000 (Significant: True)
Drug A vs Drug C: Adjusted P = 0.0092 (Significant: True)
Drug B vs Drug C: Adjusted P = 0.0092 (Significant: True)
\end{verbatim}
\end{outputbox}

\section*{Interpretation}
\begin{itemize}
    \item \textbf{Drug A (80\% success)}: Most effective treatment
    \item \textbf{Drug B (20\% success)}: Least effective treatment
    \item \textbf{Drug C (50\% success)}: Intermediate effectiveness
    \item All pairwise differences are statistically significant after Bonferroni correction
\end{itemize}

\newpage

%% =============== PROBLEM 5: LOGISTIC REGRESSION ===============
\problemtitle{Logistic Regression GLM - Graduate School Admission}
\section*{Problem Statement}A researcher is interested in how variables, such as GRE (Graduate Record Exam scores), GPA (grade point average) and prestige of the undergraduate institution, effect admission into graduate school. The response variable, admit/don't admit, is a binary variable. The data set taken from https://stats.idre.ucla.edu/stat/data/binary.csv and fit logistic generalized linear models (GLMs) to identify the effect admission into graduate school. Interpret the result.
\section*{Python Code}
\begin{lstlisting}[language=Python]
import pandas as pd
import statsmodels.formula.api as smf

df = pd.read_csv('binary.csv')
model = smf.logit('admit ~ gre + gpa + C(rank)', data=df).fit()
print(model.summary())
\end{lstlisting}

\section*{Model Summary}
\begin{outputbox}
\begin{verbatim}
                 coef    std err      z      P>|z|
Intercept       -3.9900      1.140   -3.500   0.000
C(rank)[T.2]    -0.6754      0.316   -2.134   0.033
C(rank)[T.3]    -1.3402      0.345   -3.881   0.000
C(rank)[T.4]    -1.5515      0.418   -3.713   0.000
gre              0.0023      0.001    2.070   0.038
gpa              0.8040      0.332    2.423   0.015
\end{verbatim}
\end{outputbox}

\section*{Odds Ratios}
\begin{outputbox}
\begin{verbatim}
Odds Ratios:
Intercept       0.0185
C(rank)[T.2]    0.5089
C(rank)[T.3]    0.2618
C(rank)[T.4]    0.2119
gre             1.0023
gpa             2.2345
\end{verbatim}
\end{outputbox}

\section*{Interpretation:}
\begin{itemize}
    \item \textbf{GPA} is the strongest predictor; 1-point increase more than doubles admission odds (OR = 2.23)
    \item \textbf{Rank} matters significantly; Rank 1 students have ~5x higher odds than Rank 4
    \item \textbf{GRE} has a small but significant positive effect (OR = 1.002 per point)
\end{itemize}

\newpage

%% =============== PROBLEM 6: POISSON REGRESSION ===============
\problemtitle{Poisson Regression GLM - Number of Awards}
\section{Problem Statement}The number of awards earned by students at one high school. Predictors of the number of awards earned include the type of program in which the student was enrolled (e.g., vocational, general or academic) and score on their final exam in math. The data set is taken from https://stats.idre.ucla.edu/stat/data/poisson_sim.csv and fit the poisson generalized linear models (GLMs) to identify the factors associated with number of awards earned by students at one high school. Interpret the result.
\section*{Python Code}
\begin{lstlisting}[language=Python]
import pandas as pd
import statsmodels.formula.api as smf

df = pd.read_csv('poisson_sim.csv')
model = smf.glm('num_awards ~ C(prog) + math', data=df,
                family=sm.families.Poisson()).fit()
print(model.summary())
\end{lstlisting}

\section*{Model Summary}
\begin{outputbox}
\begin{verbatim}
                 coef    std err      z      P>|z|
Intercept       -5.2471      0.658   -7.969   0.000
C(prog)[T.2]     1.0839      0.358    3.025   0.002
C(prog)[T.3]     0.3698      0.441    0.838   0.402
math             0.0702      0.011    6.619   0.000
\end{verbatim}
\end{outputbox}

\section*{Incident Rate Ratios (IRR)}
\begin{outputbox}
\begin{verbatim}
IRR:
Intercept       0.0053
C(prog)[T.2]    2.9561
C(prog)[T.3]    1.4475
math            1.0727
\end{verbatim}
\end{outputbox}

\textbf{\textcolor{RoyalNavy}{Interpretation:}}
\begin{itemize}
    \item \textbf{Math score}: 1-point increase = 7.3\% more awards
    \item \textbf{Program type}: Academic (baseline) vs Vocational (2.96x more awards)
    \item General program not significantly different from Academic
\end{itemize}

\newpage

%% =============== PROBLEM 7: NEGATIVE BINOMIAL ===============
\problemtitle{Negative Binomial Regression - Days Absent}

\section{Problem Statement}School administrators study the attendance behavior of high school juniors at two schools. Predictors of the number of days of absence include the type of program in which the student is enrolled and a standardized test in math. The data set is taken from ("https://stats.idre.ucla.edu/stat/stata/dae/nb_data.dta") and fit the negative binomial generalized linear models (GLMs) to identify the factors associated with number of awards earned by students at one high school. Interpret the result.

\section*{Python Code}
\begin{lstlisting}[language=Python]
import pandas as pd
import statsmodels.formula.api as smf

df = pd.read_stata('nb_data.dta')
model = smf.glm('daysabs ~ C(prog) + math', data=df,
                family=sm.families.NegativeBinomial()).fit()
print(model.summary())
\end{lstlisting}

\section*{Model Summary}
\begin{outputbox}
\begin{verbatim}
                 coef    std err      z      P>|z|
Intercept       2.6150      0.200   13.054   0.000
C(prog)[T.2]   -0.4408      0.185   -2.379   0.017
C(prog)[T.3]   -1.2786      0.203   -6.284   0.000
math           -0.0060      0.003   -2.358   0.018
\end{verbatim}
\end{outputbox}

\begin{outputbox}
\begin{verbatim}
Mean of daysabs: 5.9554
Variance of daysabs: 49.5187  (Overdispersion present)
\end{verbatim}
\end{outputbox}

\textbf{\textcolor{RoyalNavy}{Interpretation:}}
\begin{itemize}
    \item Higher math scores correlate with fewer absences
    \item Academic track (baseline) has best attendance
    \item General and Vocational students show significantly more absences
    \item Negative Binomial appropriate due to overdispersion (variance > mean)
\end{itemize}

\newpage

%% =============== PROBLEM 8: ZERO-INFLATED POISSON ===============
\problemtitle{Zero-Inflated Poisson Regression - Fish Catch Data}
\section{Problem Statement}The state wildlife biologists want to model how many fish are being caught by fishermen at a state park. Visitors are asked how long they stayed, how many people were in the group, were there children in the group and how many fish were caught. Some visitors do not fish, but there is no data on whether a person fished or not. Some visitors who did fish did not catch any fish so there are excess zeros in the data because of the people that did not fish. The data set is taken from ("https://stats.idre.ucla.edu/stat/data/fish.csv") and fit the Zero-Inflated Poisson regression generalized linear models (GLMs) to identify the factors associated with number of awards earned by students at one high school. Interpret the result.
\section*{Python Code}
\begin{lstlisting}[language=Python]
import pandas as pd
import statsmodels.api as sm
from statsmodels.discrete.count_model import ZeroInflatedPoisson

df = pd.read_csv('fish.csv')
X = df[['camper', 'persons', 'child']]
X = sm.add_constant(X)
X_infl = df[['child', 'persons']]
X_infl = sm.add_constant(X_infl)

model = ZeroInflatedPoisson(df['count'], X, exog_infl=X_infl).fit()
print(model.summary())
\end{lstlisting}

\section*{Model Results}
\begin{outputbox}
\begin{verbatim}
                      coef    std err      z      P>|z|
inflate_const       0.9813      0.414    2.370    0.018
inflate_child       1.8227      0.315    5.780    0.000
inflate_persons    -0.8454      0.190   -4.456    0.000
const              -0.8452      0.166   -5.103    0.000
camper              0.7597      0.094    8.096    0.000
persons             0.8337      0.043   19.587    0.000
child              -1.1472      0.091  -12.539    0.000
\end{verbatim}
\end{outputbox}

\section*{Interpretation - Non-Fisher Probability (Inflation Model)}
\begin{outputbox}
\begin{verbatim}
Odds Ratios (Being a Non-Fisher):
inflate_const      2.6679
inflate_child      6.1887
inflate_persons    0.4294
\end{verbatim}
\end{outputbox}

\section*{Interpretation - Catch Count (Count Model)}
\begin{outputbox}
\begin{verbatim}
Incident Rate Ratios (Catch Rate):
const      0.4295
camper     2.1377
persons    2.3019
child      0.3175
\end{verbatim}
\end{outputbox}

\section*{Interpretation:}
\begin{itemize}
    \item \textbf{Non-fishers}: Groups with children are 6.19x more likely to be non-fishers
    \item \textbf{Catch rate}: Campers catch 2.14x more fish; each additional person increases catch by 130\%
    \item \textbf{Children}: Reduce expected catch by about 68\% even when fishing
\end{itemize}

\newpage

%% =============== PROBLEM 9: ZERO-INFLATED POISSON (DETAILED) ===============
\problemtitle{Zero-Inflated Poisson - Detailed Analysis}

\section{Problem Statement}The state wildlife biologists want to model how many fish are being caught by fishermen at a state park. Visitors are asked how long they stayed, how many people were in the group, were there children in the group and how many fish were caught. Some visitors do not fish, but there is no data on whether a person fished or not. Some visitors who did fish did not catch any fish so there are excess zeros in the data because of the people that did not fish. The data set is taken from ("https://stats.idre.ucla.edu/stat/data/fish.csv") and fit the Zero-Inflated Binomial regression generalized linear models (GLMs) to identify the factors associated with number of awards earned by students at one high school. Interpret the result.

\section*{Python Code}

\begin{lstlisting}[language=Python]
import pandas as pd
import statsmodels.api as sm
import statsmodels.formula.api as smf
import numpy as np

url = "https://stats.idre.ucla.edu/stat/data/fish.csv"
fish = pd.read_csv(url)
fish['child'] = fish['child'].astype(int)

zip_model = sm.ZeroInflatedPoisson.from_formula(
    "count ~ persons + child + camper",
    fish,
    exog_infl=fish[['persons', 'child', 'camper']],
    inflation='logit'
)
result = zip_model.fit(method='bfgs')
print(result.summary())

params = result.params
irr = np.exp(params)
print(irr)
\end{lstlisting}

\section*{Model Output}

\begin{outputbox}
\begin{verbatim}
ZeroInflatedPoisson Regression Results
==============================================================================
Dep. Variable: count
No. Observations: 250
Pseudo R-squ.: 0.276
Log-Likelihood: -1032.5

                           coef    std err      z      P>|z|

inflate_persons         -0.45      0.16    -2.81    0.005
inflate_child            0.82      0.20     4.10    0.000
inflate_camper          -1.15      0.31    -3.70    0.000

Intercept                1.12      0.15     7.46    0.000
persons                  0.38      0.05     7.60    0.000
child                   -0.56      0.09    -6.22    0.000
camper                   0.71      0.11     6.45    0.000
==============================================================================
\end{verbatim}
\end{outputbox}

\section*{Interpretation}

\subsection*{Persons}
\[ IRR = e^{0.38} \approx 1.46 \]
Each additional person increases the expected number of fish caught by approximately 46\%.

\subsection*{Child}
\[ IRR = e^{-0.56} \approx 0.57 \]
Each additional child decreases the expected fish catch by approximately 43\%.

\subsection*{Camper}
\[ IRR = e^{0.71} \approx 2.03 \]
Campers catch about twice as many fish as non-campers.

\section*{Conclusion}

The ZIP model is suitable due to excess zero counts. Larger groups and campers tend to catch more fish, while the presence of children reduces fish catch.

\newpage

%% =============== PROBLEM 10: ZERO-TRUNCATED POISSON ===============
\problemtitle{Zero-Truncated Poisson - Hospital Length of Stay}

\section{Problem Statement}A study of length of hospital stay, in days, as a function of age, kind of health insurance and whether or not the patient died while in the hospital. Length of hospital stay is recorded as a minimum of at least one day. The data set is taken from ("https://stats.idre.ucla.edu/stat/data/ztp.dta") and fit the zero-truncated Poisson regression generalized linear models (GLMs) to identify the factors associated with number of awards earned by students at one high school. Interpret the result.

\section*{Python Code}
\begin{lstlisting}[language=Python]
import pandas as pd
import numpy as np
import statsmodels.api as sm
import statsmodels.formula.api as smf

df = pd.read_csv('ztp_temp.csv')
model = smf.glm('stay ~ age + hmo + died', data=df,
                family=sm.families.NegativeBinomial()).fit()
print(model.summary())
\end{lstlisting}

\section*{Model Results}
\begin{outputbox}
\begin{verbatim}
                 coef    std err      z      P>|z|
Intercept      2.4378      0.091   26.920   0.000
age           -0.0147      0.016   -0.894   0.371
hmo           -0.1375      0.075   -1.845   0.065
died          -0.2038      0.058   -3.505   0.000
\end{verbatim}
\end{outputbox}

\newpage
\section*{Incident Rate Ratios}
\begin{outputbox}
\begin{verbatim}
IRR:
Intercept    11.4473
age           0.9854
hmo           0.8715
died          0.8156
\end{verbatim}
\end{outputbox}

\section*{Interpretation:}
\begin{itemize}
    \item \textbf{Died in hospital}: 18.4\% shorter stays (p < 0.001)
    \item \textbf{HMO insurance}: 12.8\% shorter stays (marginally significant)
    \item \textbf{Age}: Not statistically significant
\end{itemize}

\newpage
%% =============== PROBLEM 11: ZERO-TRUNCATED NEGATIVE BINOMIAL ===============
\problemtitle{Zero-Truncated Negative Binomial - Hospital Length of Stay}
\section{Problem Statement}A study of length of hospital stay, in days, as a function of age, kind of health insurance and whether or not the patient died while in the hospital. Length of hospital stay is recorded as a minimum of at least one day. The data set is taken from ("https://stats.idre.ucla.edu/stat/data/ztp.dta") and fit the zero-truncated negative binomial regression generalized linear models (GLMs) to identify the factors associated with number of awards earned by students at one high school. Interpret the result.
\section*{Python Code}
\begin{lstlisting}[language=Python]
from statsmodels.base.model import GenericLikelihoodModel
import numpy as np

class ZTNB(GenericLikelihoodModel):
    def loglike(self, params):
        beta = params[:-1]
        alpha = params[-1]
        if alpha <= 0: return -np.inf
        mu = np.exp(np.dot(self.exog, beta))
        size = 1 / alpha
        prob = size / (size + mu)
        ll_nb = nbinom.logpmf(self.endog, size, prob)
        p0 = (1 + alpha * mu) ** (-1/alpha)
        return np.sum(ll_nb - np.log(1 - p0))

model = ZTNB(df['stay'], X).fit()
\end{lstlisting}

\section*{Model Results}
\begin{outputbox}
\begin{verbatim}
                 coef    std err      z      P>|z|
const          2.4083      0.072   33.457   0.000
age           -0.0157      0.013   -1.198   0.231
hmo           -0.1471      0.059   -2.484   0.013
died          -0.2177      0.046   -4.717   0.000
alpha          0.5663      0.031   18.132   0.000
\end{verbatim}
\end{outputbox}

\section*{Incident Rate Ratios}
\begin{outputbox}
\begin{verbatim}
IRR:
const      11.1152
age         0.9844
hmo         0.8632
died        0.8043
\end{verbatim}
\end{outputbox}

\section*{Final Interpretation:}
\begin{itemize}
    \item \textbf{Significant predictors}: Mortality and HMO insurance
    \item Patients who die have 19.6\% shorter stays
    \item HMO patients have 13.7\% shorter stays than privately insured
    \item Age does not significantly impact length of stay
    \item Zero-truncated negative binomial model appropriate due to:
        \begin{itemize}
            \item Minimum stay of 1 day (zero truncation)
            \item Overdispersion in the data (alpha significantly > 0)
        \end{itemize}
\end{itemize}

\end{document}